
This work demonstrates that benchmark popularity correlates with larger generalization failures in machine learning models. Models trained on the high-popularity dataset CIFAR-10 exhibit a generalization gap of 18.86 percentage points when evaluated on CINIC-10, while models trained on the lower-popularity SVHN dataset show a gap of -2.35 percentage points to SVHN-Extra---a difference of 21.21 percentage points despite identical architectures and training protocols. A bibliometric analysis finds that high-use datasets receive 7.56 times more optimization-focused papers than low-use datasets (p=0.027), consistent with a benchmark co-evolution hypothesis. However, testing texture bias as the proposed mechanism, experiments find that ResNet-18 exhibits lower texture bias (0.126) than VGG-11 (0.161), ruling out texture exploitation as the causal pathway. These results suggest that benchmark selection itself may confound model evaluation. The findings support complementing popular benchmarks with diverse, less-optimized alternatives for deployment validation.

